% ============================================================
% Section 3 — Structural model setup
% Bridges CWS dynamic programming problem to Hotz--Miller log-odds CCP
% Author: EK   Date: 2026-05-20
% Requires: amsmath, amssymb (and \citep from natbib/biblatex)
% ============================================================

\section{A Dynamic Model of CWS Compliance}
\label{sec:model}

\subsection{Environment}

Time is discrete and indexed by $t = 1, 2, \ldots$. A community water system
(CWS) $i$ chooses an action $a_{it} \in \{0, 1\}$ each period, where $a_{it} = 1$
denotes \emph{compliance} (conducting tests and reporting them) and $a_{it} = 0$
denotes \emph{negligence}, which produces a monitoring-and-reporting (MR)
violation. Compliance carries a per-period administrative cost $k_{\text{comply}}$;
negligence avoids this cost but raises the probability of future enforcement and
carries an implicit reputational/operational cost $k_{\text{MR}}$.

\paragraph{State space.} The CWS faces a state $s_{it} = (x_{it}, \varepsilon_{it})$
with two components:
\begin{itemize}
\item $x_{it} \in \mathcal{X}$ is the \emph{observed} state, including
      enforcement status ($0 = $ no active enforcement, $1 = $ currently under
      enforcement) and a discretized mining-exposure measure. $x_{it}$ is
      observed by both the CWS and the econometrician.
\item $\varepsilon_{it} = \big(\varepsilon_{it}(0), \varepsilon_{it}(1)\big)
      \in \mathbb{R}^2$ is a vector of action-specific payoff shocks. The CWS
      observes $\varepsilon_{it}$ before choosing $a_{it}$; the econometrician
      never observes it. These shocks capture transitory variation in the
      relative attractiveness of complying versus violating (e.g., temporary
      staffing constraints, lab availability, managerial attention).
\end{itemize}

Per-period payoff is additively separable in $\varepsilon$:
\begin{equation}
U(a, x_{it}) + \varepsilon_{it}(a),
\label{eq:flow}
\end{equation}
where $U(1, x) = -k_{\text{comply}} - c(x)$ and $U(0, x) = -k_{\text{MR}} - c(x)$,
and $c(x)$ is the direct cost of being in enforcement state $x$ (normalized so
that $c(0) = 0$ and $c(1) = 1$).\footnote{Because $c(x)$ depends only on the
enforcement state and not on the action, it cancels from the flow-utility
difference: $U(0, x) - U(1, x) = k_{\text{comply}} - k_{\text{MR}}$,
independent of $x$. The enforcement cost $c(x)$ enters the estimating equation
only through the continuation value $\bar V(x)$ in \eqref{eq:Vbar}, and only
through the difference $\bar V(1) - \bar V(0)$.}

\paragraph{Distributional and independence assumptions.} Following
\citet{rust1987optimal} and \citet{hotz1993conditional}, we impose:
\begin{description}
\item[(A1) Additive separability.] Payoffs take the form in
      \eqref{eq:flow}.
\item[(A2) T1EV shocks.] The $\varepsilon_{it}(a)$ are i.i.d.\ across $i$,
      $t$, and $a$, drawn from a Type-1 Extreme Value distribution. The scale
      is normalized to one ($\sigma_\varepsilon \equiv 1$); together with
      $c(1) = 1$ this fixes the model's units, so $k_{\text{MR}}$ is identified
      in units of the enforcement-state cost $c(1)$ rather than dollars.
\item[(A3) Conditional independence.] The transition density factors as
      $f(x', \varepsilon' \mid x, \varepsilon, a) = f_x(x' \mid x, a)
      \cdot f_\varepsilon(\varepsilon')$. That is, conditional on the current
      observed state and action, next period's observed state $x'$ does not
      depend on the current $\varepsilon$, and next period's $\varepsilon'$ is
      drawn independently of the past.
\end{description}
Assumption (A3) is the Conditional Independence Assumption of
\citet{rust1987optimal}. It is what makes $\varepsilon$ a transitory taste
shock rather than a persistent unobservable, and is essential for the
representation derived below.

\subsection{The CWS's dynamic programming problem}

The CWS chooses a policy $\{a_{it}\}$ to maximize the expected
discounted sum of payoffs, $\mathbb{E}\sum_{j=0}^{\infty} \beta^{j}
\big[ U(a_{i,t+j}, x_{i,t+j}) + \varepsilon_{i,t+j}(a_{i,t+j}) \big]$,
with discount factor $\beta \in (0, 1)$ (set to $\beta = 0.95$ in
estimation). By the principle of optimality, the value function over the
\emph{full} state satisfies the Bellman equation
\begin{equation}
V(x, \varepsilon) = \max_{a \in \{0, 1\}} \Big\{
U(a, x) + \varepsilon(a) + \beta \, \mathbb{E}\!\left[ V(x', \varepsilon')
\,\big|\, x, a \right] \Big\}.
\label{eq:bellman_full}
\end{equation}

\subsection{Integrating out the private shock}

Equation~\eqref{eq:bellman_full} is intractable as written because $V$ depends on
the unobserved $\varepsilon$. Under (A3), however, we can integrate
$\varepsilon'$ out of the continuation. Define the \emph{choice-specific value
function}
\begin{equation}
\bar v(a, x) \;\equiv\; U(a, x) + \beta \sum_{x'} \bar V(x') \, P(x' \mid x, a),
\label{eq:vbar}
\end{equation}
where $P(x' \mid x, a) \equiv f_x(x' \mid x, a)$ is the action-conditional
transition probability between observed states, and $\bar V(x)$ is the
\emph{ex-ante} (or \emph{integrated}) value function
\begin{equation}
\bar V(x) \;\equiv\; \mathbb{E}_\varepsilon \!\left[ \max_{a}
\big\{ \bar v(a, x) + \varepsilon(a) \big\} \right].
\label{eq:Vbar}
\end{equation}
Under (A2), equation \eqref{eq:Vbar} has a closed form \citep{train2009discrete}:
\begin{equation}
\bar V(x) = \log \!\sum_{a \in \{0,1\}} \exp\!\left( \bar v(a, x) \right) + C,
\label{eq:logsumexp}
\end{equation}
where $C$ is a constant (the Euler--Mascheroni term in the exact McFadden
formula) that is the same for every $x$ and therefore cancels in every
difference $\bar V(x) - \bar V(x')$ taken below.
Equations~\eqref{eq:vbar}--\eqref{eq:logsumexp} together constitute a fixed
point in $\bar v$ and $\bar V$ that depends only on the observed state $x$.

\subsection{The Hotz--Miller representation}

Given the choice-specific values $\bar v(a, x)$, the optimal action solves a
static problem with deterministic payoffs and T1EV shocks:
\begin{equation*}
a^*(x, \varepsilon) = \arg\max_{a \in \{0,1\}}
\big\{ \bar v(a, x) + \varepsilon(a) \big\}.
\end{equation*}
Standard logit algebra \citep[Ch.~3]{train2009discrete} yields the conditional
choice probability
\begin{equation}
\Pr(a \mid x) = \frac{\exp\!\left( \bar v(a, x) \right)}
{\sum_{a' \in \{0, 1\}} \exp\!\left( \bar v(a', x) \right)}.
\label{eq:ccp}
\end{equation}
Taking log-odds between negligence ($a = 0$, MR) and compliance ($a = 1$):
\begin{equation}
\boxed{\;
\log \!\left[ \frac{\Pr(\text{MR} \mid x)}{\Pr(\text{comply} \mid x)} \right]
\;=\; \bar v(0, x) - \bar v(1, x)
\;\equiv\; \Delta \bar v(x).
\;}
\label{eq:logodds}
\end{equation}
Equation~\eqref{eq:logodds} is the \citet{hotz1993conditional} representation
theorem applied to our setting: the left-hand side is observed in the data
(empirical log-odds within state $x$), and the right-hand side is a closed-form
expression in structural primitives.

\subsection{Computing the integrated value function in practice}
\label{sec:vbar_computation}

Equation~\eqref{eq:logsumexp} defines $\bar V(x)$ as the fixed point of a
nonlinear logsumexp system. Solving it directly would require the nested
fixed-point algorithm of \citet{rust1987optimal}, which the
\citet{hotz1993conditional} representation is designed to avoid. We instead
use the \citet{arcidiaconomiller2011} inversion, which expresses $\bar V$ as a
closed-form linear function of observed CCPs.

\paragraph{The Arcidiacono--Miller inversion.}
Combining the CCP formula \eqref{eq:ccp} with the logsumexp expression
\eqref{eq:logsumexp} and taking logs yields the identity\footnote{Algebra:
\eqref{eq:ccp} gives $\bar v(a, x) =
\log \hat\Pr(a \mid x) + \log \sum_{a'} \exp(\bar v(a', x))$;
substituting \eqref{eq:logsumexp} delivers \eqref{eq:am_inversion}.}
\begin{equation}
\bar V(x) = \bar v(a, x) - \log \hat\Pr(a \mid x) + C,
\qquad \text{for any } a \in \{0, 1\}.
\label{eq:am_inversion}
\end{equation}
Equation~\eqref{eq:am_inversion} is exact under T1EV — it inverts the CCP
mapping to express the integrated value in terms of any single choice-specific
value plus a CCP correction. Choosing the comply action ($a = 1$) as the
reference and substituting \eqref{eq:am_inversion} into the Bellman
\eqref{eq:vbar} for $a = 1$ gives a linear system in $\bar v(1, \cdot)$ alone:
\begin{equation}
\bar v(1, x) = U(1, x) + \beta \sum_{x'} P(x' \mid x, 1)
\big[ \bar v(1, x') - \log \hat\Pr(1 \mid x') \big]
+ \text{const}.
\label{eq:vbar1_system}
\end{equation}
Stacking over the two enforcement states $x \in \{0, 1\}$, let
$\boldsymbol{\bar v}_1 = (\bar v(1, 0), \bar v(1, 1))'$,
$\boldsymbol{U}_1 = (U(1, 0), U(1, 1))'$, $P_1$ the $2 \times 2$ comply-action
transition matrix, and $\boldsymbol{\ell}_1 = (\log \hat\Pr(1 \mid 0),
\log \hat\Pr(1 \mid 1))'$. Then
\begin{equation}
\boldsymbol{\bar v}_1 = (I - \beta P_1)^{-1}
\left[ \boldsymbol{U}_1 - \beta P_1 \boldsymbol{\ell}_1 \right]
+ \text{const} \cdot \mathbf{1},
\label{eq:vbar1_solve}
\end{equation}
where the additive constant collects terms that are equal across $x$. Equation~\eqref{eq:vbar1_solve} is a $2 \times 2$ linear
solve using \emph{only} the comply-action transition matrix $P_1$, the
flow-utility primitives $U(1, \cdot)$, and the observed comply-CCPs
$\hat\Pr(1 \mid x)$.

\paragraph{What enters the estimating equation.} Applying
\eqref{eq:am_inversion} to both states and differencing:
\begin{equation}
\bar V(1) - \bar V(0) = \big[\bar v(1, 1) - \bar v(1, 0)\big]
- \big[\log \hat\Pr(1 \mid 1) - \log \hat\Pr(1 \mid 0)\big].
\label{eq:vbar_diff}
\end{equation}
The constant in \eqref{eq:vbar1_solve} cancels in the difference. The first
bracket comes from \eqref{eq:vbar1_solve}; the second is computed directly
from sample CCPs. Together \eqref{eq:vbar1_solve}--\eqref{eq:vbar_diff} replace
a nonlinear Bellman fixed point with one linear system and a log-CCP
difference.

\paragraph{Simplification: $k_{\text{comply}}$ drops out of $\bar V(1) - \bar V(0)$.}
The flow utility under comply is $U(1, x) = -k_{\text{comply}} - c(x)$.
Decompose $\boldsymbol{U}_1 = -k_{\text{comply}} \cdot \mathbf{1} -
\boldsymbol{c}$, where $\boldsymbol{c} = (c(0), c(1))' = (0, 1)'$. Because
$P_1$ is row-stochastic, $P_1 \mathbf{1} = \mathbf{1}$ and hence
$(I - \beta P_1)^{-1} \mathbf{1} = (1 - \beta)^{-1} \mathbf{1}$, so the
$k_{\text{comply}}$ term contributes $-k_{\text{comply}} / (1 - \beta)$
\emph{uniformly} to both entries of $\boldsymbol{\bar v}_1$. It therefore
cancels in the difference $\bar v(1, 1) - \bar v(1, 0)$. Substituting into
\eqref{eq:vbar_diff} yields a closed form for the continuation-value
difference that does \emph{not} depend on $k_{\text{comply}}$:
\begin{equation}
\bar V(1) - \bar V(0) =
- \Delta_x \!\big[(I - \beta P_1)^{-1} \boldsymbol{c}\big]
- \beta \, \Delta_x \!\big[(I - \beta P_1)^{-1} P_1 \boldsymbol{\ell}_1\big]
- \Delta_x \boldsymbol{\ell}_1,
\label{eq:vbar_diff_final}
\end{equation}
where $\Delta_x [w] \equiv w(1) - w(0)$ denotes the second-minus-first entry
of a 2-vector indexed by the enforcement state $x \in \{0, 1\}$. The three
terms have clean interpretations: the first is the discounted enforcement
burden gap under the comply policy; the second is a T1EV-entropy correction
weighted by the comply-action transitions; the third is the contemporaneous
log-CCP gap.

This simplification matters operationally. It means $\bar V(1) - \bar V(0)$
can be constructed from data alone — the action-conditional transition matrix
$P_1$, the empirical comply-CCPs $\hat\Pr(1 \mid x)$, and the normalized
enforcement cost vector $\boldsymbol{c}$ — with no need to iterate over the
unknown parameter $k_{\text{comply}}$. The
intercept of the estimating equation \eqref{eq:estimating} then identifies the
\emph{difference} $k_{\text{comply}} - k_{\text{MR}}$ directly, consistent with
the fact that the additive level of $k_{\text{comply}}$ is not separately
identified from the data.\footnote{Equations \eqref{eq:vbar1_solve} and \eqref{eq:vbar_diff}
are exact at population CCPs. With sample CCPs $\hat\Pr(a \mid x)$ they
inherit $O_p(N_x^{-1/2})$ sampling error, where $N_x$ is the cell size
($\approx 780$ in our application). This propagates into the regression
slope through $\bar V(1) - \bar V(0)$ as a generated-regressor error;
standard errors in Section~\ref{sec:estimation} account for it via the
bootstrap.}

\paragraph{Comparison with the naive forward-simulation approach.}
A frequent shortcut in the applied literature is to compute the value of the
\emph{empirically observed policy}
$V^*(x) = (I - \beta P^*)^{-1} c$, where $P^*(x' \mid x) = \sum_a
\hat\Pr(a \mid x) P(x' \mid x, a)$, and treat $V^*(1) - V^*(0)$ as a proxy for
$\bar V(1) - \bar V(0)$. This is *not* exact: $V^*$ ignores both the flow
utility of choices under the policy and the T1EV entropy term
$-\sum_a \hat\Pr(a \mid x') \log \hat\Pr(a \mid x')$, which
varies with $x'$ through the CCPs. The difference $\bar V(1) - \bar V(0) -
[V^*(1) - V^*(0)]$ is in general nonzero and state-dependent, biasing the
regression slope. We use \eqref{eq:vbar1_solve}--\eqref{eq:vbar_diff}
throughout and report $V^*$-based estimates in
Appendix~\ref{app:naive_forwardsim} only as a diagnostic.

\paragraph{No estimation fixed point.} The construction
\eqref{eq:vbar1_solve}--\eqref{eq:vbar_diff_final} requires no iteration.
Action-conditional transition probabilities $P(x' \mid x, a)$ and CCPs
$\hat\Pr(a \mid x)$ are both estimated by tabulation; $\bar V(1) - \bar V(0)$
is a closed-form function of them; the structural
parameter $k_{\text{MR}}$ then comes from a single linear regression
\eqref{eq:estimating}. There is no nested fixed point as in \citet{rust1987optimal}
and no inner loop on $k_{\text{comply}}$ (which cancels, as shown above).

\paragraph{Caveat: optimality of the empirical policy.} The inversion
\eqref{eq:am_inversion} does, however, require that the empirical
$\hat\Pr(a \mid x)$ be generated by the structural CCP \eqref{eq:ccp} — i.e.,
that CWSs play the policy implied by the parameters we are recovering. This
is an \emph{equilibrium} assumption (the data come from a stationary optimal
policy), distinct from the \emph{computational} fixed point that the AM
inversion replaces. It is the standard two-step assumption in the CCP
literature \citep{hotzmillersanders1994, aguirregabiriamira2002,
bajaribenkardlevin2007} and would be violated by, e.g., systematic learning or
out-of-equilibrium play during the sample period. We report estimates from one
iteration of the recursive updating procedure of \citet{aguirregabiriamira2002}
as a robustness check in Section~\ref{sec:robustness_ccp_iter}.

\subsection{Decomposition and estimating equation}

Substituting \eqref{eq:vbar} into \eqref{eq:logodds} and collecting terms:
\begin{equation}
\Delta \bar v(x) = \underbrace{\big[ U(0, x) - U(1, x) \big]}_{\text{flow utility difference}}
+ \beta \underbrace{\sum_{x'} \bar V(x')
\big[ P(x' \mid x, 0) - P(x' \mid x, 1) \big]}_{\text{continuation value difference}}.
\label{eq:decomp}
\end{equation}
With two enforcement states ($x' \in \{0, 1\}$), the continuation term reduces
to
\begin{equation*}
\beta \cdot \big[\bar V(1) - \bar V(0)\big] \cdot \Delta \pi(x),
\quad \text{where} \quad
\Delta \pi(x) \equiv \Pr(x' = 1 \mid x, \text{MR}) - \Pr(x' = 1 \mid x, \text{comply}).
\end{equation*}
$\Delta \pi(x) > 0$ encodes the dynamic deterrence channel: negligence raises
the probability of being under enforcement next period. The estimating
equation is therefore
\begin{equation}
\log \!\left[ \frac{\Pr(\text{MR} \mid x)}{\Pr(\text{comply} \mid x)} \right]
= \underbrace{(k_{\text{comply}} - k_{\text{MR}})}_{\text{intercept}}
+ \underbrace{\beta \big[\bar V(1) - \bar V(0)\big]}_{\text{slope}}
\cdot \Delta \pi(x) + \gamma \, \hat{v}_t,
\label{eq:estimating}
\end{equation}
where $\hat{v}_t$ is the first-stage residual from regressing mining exposure on
the ARP $\times$ sulfur instrument, included as a control function
\citep{petrintrain2010, wooldridge2015control}. Under the scale normalization
$\sigma_\varepsilon = 1$ from (A2), the intercept identifies
$k_{\text{comply}} - k_{\text{MR}}$ directly once the slope identifies
$\bar V(1) - \bar V(0)$; with $c(1) = 1$ this is in units of the
enforcement-state cost.

\subsection{Role of the private shock $\varepsilon$}

The private shock $\varepsilon_{it}$ plays three distinct roles that are worth
distinguishing from \emph{permanent} unobserved heterogeneity (e.g.,
``compliant-by-culture'' types) in the sense of \citet{arcidiaconomiller2011}.
First, $\varepsilon$ rationalizes within-state variation in observed choices:
without it, all CWSs sharing the same $x$ would make identical decisions and
choice \emph{probabilities} would be degenerate. Second, the T1EV assumption
(A2) yields the closed-form invertible mapping between CCPs and
choice-specific value differences in~\eqref{eq:logodds}, allowing us to
estimate the model without solving the full Bellman fixed point. Third, the
T1EV scale is the only free scale in the model, which is why we normalize it
to one in (A2); together with $c(1) = 1$ this fixes the units in which
$k_{\text{MR}}$ is recovered. Crucially,
(A3) requires that $\varepsilon$ be i.i.d.\ across $(i, t, a)$, so it is a
transitory taste shock, not a persistent type. We discuss permanent unobserved
heterogeneity as a robustness check in
Section~\ref{sec:robustness_unobs_heterogeneity}.
